\documentclass[12pt]{article}
\usepackage{fullpage}
\usepackage{amsmath,amssymb,amsthm}

\newcommand{\cL}{\mathcal L}
\newcommand{\Dr}{\operatorname{Dr}}
\newcommand{\Trop}{\operatorname{Trop}}
\newcommand{\rk}{\operatorname{rk}}
\newcommand{\cl}{\operatorname{cl}}
\newcommand{\op}{\operatorname{op}}

\newtheorem{theorem}{Theorem}
\newtheorem{lemma}{Lemma}
\newtheorem{proposition}{Proposition}

\begin{document}

\begin{center}
{\bf A canonical opposition involution on the relative Dressian}
\end{center}

\section*{Conventions and local flag input}

Let $E$ be the ground set, let $r=\rk M$, and let $\mu_M$ be the
trivial rank $r$ valuated matroid supported on the bases of $M$.  We
use the standard valuated-matroid model of the relative Dressian: a
rank $k$ point is a projective rank $k$ valuated matroid
$\nu$ with $\nu\preceq\mu_M$, where $\preceq$ denotes valuated quotient.
Brandt--Eur--Zhang identify this with inclusion of projective tropical
linear spaces:
\[
        \Trop(\theta)\subseteq \Trop(\nu)
        \quad\Longleftrightarrow\quad \theta\preceq\nu .
\]
For valuated matroids $p,q$ of ranks $a\le b$, the quotient relation is
the full incidence-Plucker condition
\[
 \min_{i\in J\setminus I}\{p(I\cup i)+q(J\setminus i)\}
 \quad\hbox{is attained at least twice}                     \tag{A}
\]
for all $I\in\binom E{a-1}$ and $J\in\binom E{b+1}$, together with the
Plucker relations for $p$ and $q$.  Coordinates equal to $\infty$ are
allowed; an all-$\infty$ minimum is attained at every index.

The only external cryptomorphism used below is the following local form.
It is the $0$-$1$ specialization of Murota--Shioura's local exchange
criterion for $M^\natural$-concave functions on generalized
polymatroids (Theorem 4.2 of their simpler-exchange paper, equivalent
to Theorem 6.15 of Murota's book), and the tropical perfect-tract
flag-matroid cryptomorphism of Jarra--Lorscheid (Theorem C and
Proposition 2.21).  The hypotheses explicitly allow non-uniform
ordinary supports and $\infty$-coordinates: the effective domain is the
ordinary flag matroid support, an integral generalized polymatroid, and
the displayed tropical bend relations are the local exchange axioms on
that domain.

\begin{lemma}[local flag criterion]\label{lem:localflag}
Let $p,q$ be valuated matroids of ranks $a\le b$ whose ordinary supports
are in quotient order.  Then the full relations {\rm (A)} are equivalent
to the nested relations
\[
 \min_{i\in T\setminus S}\{p(S\cup i)+q(T\setminus i)\}
 \quad\hbox{attained at least twice}                         \tag{N}
\]
for all $S\subseteq T$, $|S|=a-1$, $|T|=b+1$.

Moreover, let $g_0,\ldots,g_r$ be nonzero homogeneous tropical
coordinate vectors whose ordinary supports form a complete ordinary flag
matroid with the trivial rank $0$ support and top support $M$.  Then the
sequence is a complete valuated flag matroid if and only if every
adjacent pair satisfies the local three-term relations
\[
 \min\{g_j(Sx)+g_{j+1}(Syz),\,
        g_j(Sy)+g_{j+1}(Sxz),\,
        g_j(Sz)+g_{j+1}(Sxy)\}
\]
for all $|S|=j-1$ and distinct $x,y,z\notin S$.  Finally, if
$\nu\preceq\mu_M$ has rank $k$, then it can be inserted in such a
complete flag by
\[
 h_j(B)=
 \begin{cases}
 \min\{\nu(C): B\subseteq C,\ |C|=k\}, & j\le k,\ |B|=j,\\
 \min\{\nu(C): C\subseteq B,\ |C|=k\}, & j\ge k,\ B\in\mathcal I_j(M),\\
 \infty, & j\ge k,\ B\notin\mathcal I_j(M).
 \end{cases}                                                \tag{1}
\]
\end{lemma}

\begin{proof}
For two ranks, the support condition says exactly that the union of the
two basis layers is an ordinary flag matroid, i.e. a $0$-$1$
generalized polymatroid.  On this domain the local exchange axiom of
Murota--Shioura is precisely (N), while the global exchange axiom is the
basis-exchange form of valuated quotients, equivalently the arbitrary
incidence relations (A) by Brandt--Eur--Zhang, Proposition 4.2.3.  This
proves the first assertion, including the nonuniform and $\infty$ cases.

For a complete flag, Jarra--Lorscheid's perfect-tract theorem says that
the adjacent Plucker flag relations for the tropical tract are
equivalent to a flag $\mathbb T$-matroid; under their Proposition 2.21
this is the same object as a valuated flag matroid, hence as a chain of
tropical linear spaces.  The displayed adjacent relations are exactly
those relations in ranks $j,j+1$.  Formula (1) is the
Dress--Terhalle--Murota well-layering of the independent-set valuation
associated with $\nu$: the layers $j\le k$ are valuated truncations and
the layers $j\ge k$ are relative Higgs elongations inside $M$.  The top
layer is projectively $\mu_M$: if $B'=B-e+f$ are adjacent $M$-bases and
$C\subseteq B$ minimizes the second line of (1), then either $e\notin C$
and $C\subseteq B'$, or the incidence relation for
$\nu\preceq\mu_M$ with lower set $C-e$ and upper set $B\cup f$ produces
a $k$-subset of $B'$ of value at most $\nu(C)$; symmetry gives equality
along the basis graph.
\end{proof}

\section*{1. Consequences of the flat-lattice duality}

\begin{lemma}\label{lem:modular}
If $L=\cL(M)$ admits an order-reversing involution
$x\mapsto x^\perp$, then $L$ is modular and
$\rk(x^\perp)=r-\rk(x)$.
\end{lemma}

\begin{proof}
The anti-automorphism sends $\hat0$ to $\hat1$ and reverses maximal
chains, hence sends rank $i$ flats to rank $r-i$ flats.  The flat
lattice is semimodular.  Applying the anti-automorphism to the
semimodular inequality gives the reverse inequality, so equality holds
for every pair of flats.  Thus $L$ is modular.
\end{proof}

Thus $(x\vee y)^\perp=x^\perp\wedge y^\perp$ and
$(x\wedge y)^\perp=x^\perp\vee y^\perp$.

\begin{lemma}[flat constancy]\label{lem:flatconst}
Let $\nu\preceq \mu_M$ have rank $k$.  If $B,B'$ are independent
$k$-subsets of $M$ with $\cl(B)=\cl(B')$, then $\nu(B)=\nu(B')$.
Also $\nu(A)=\infty$ whenever $A$ is dependent in $M$.
\end{lemma}

\begin{proof}
The support matroid of $\nu$ is an ordinary quotient of $M$, so each of
its bases is independent in $M$.  The case $k=0$ is trivial.  For
$k>0$, the basis graph of the restriction of $M$ to a rank $k$ flat is
connected, so it suffices to compare $B=A\cup b$ and $B'=A\cup b'$
with $|A|=k-1$ and $\cl(Ab)=\cl(Ab')$.  Choose $C$ of size $r-k$ such
that $AbC$ is an $M$-basis; then $Ab'C$ is also an $M$-basis.  In the
nested incidence relation for $\nu\preceq\mu_M$ with lower set $A$ and
upper set $A\cup\{b,b'\}\cup C$, the only finite $\mu_M$-terms are
those indexed by $b$ and $b'$, since deleting an element of $C$ leaves
the dependent set $Abb'$.  Therefore $\nu(Ab)=\nu(Ab')$.
\end{proof}

We write $\bar\nu(X)=\nu(B)$ for a rank $k$ flat $X=\cl(B)$.

\section*{2. The opposition transform and supports}

For a rank $k$ quotient $\nu\preceq\mu_M$ define a rank $r-k$ vector by
\[
 \nu^{\perp_M}(A)=
 \begin{cases}
   \bar\nu\bigl((\cl A)^\perp\bigr),
      & A\text{ independent in }M,\ |A|=r-k,\\
   \infty,&\text{otherwise.}
 \end{cases}                                                 \tag{2}
\]
This is compatible with projective rescaling, and applying (2) twice
recovers the original flat coordinates.

\begin{lemma}[ordinary polar rank]\label{lem:ordinary}
Let $N$ be an ordinary quotient of $M$ of rank $k$, with rank function
$\rho_N$.  The sets
\[
 {\mathcal I}^{\perp}=\{A:\ A\text{ is }M\text{-independent and }
       \rho_N((\cl A)^\perp)=k\}
\]
are the independent sets of a matroid $N^{\perp_M}$.  For every flat
$X$ of $M$,
\[
        \rho_{N^{\perp_M}}(X)=\rk_M(X)-k+\rho_N(X^\perp).       \tag{3}
\]
Consequently, if $N_p$ is a quotient of $N_q$, then
$N_q^{\perp_M}$ is a quotient of $N_p^{\perp_M}$.
\end{lemma}

\begin{proof}
The right side of (3), extended from subsets to their $M$-closures, is
a matroid rank function.  It is zero at $\hat0$, is monotone because a
quotient rank can increase along an interval no faster than the
$M$-rank, and is submodular because the $M$-rank is modular and
$\rho_N$ is submodular, while $^\perp$ interchanges meet and join.
For an $M$-independent set $A$ with $Y=\cl(A)$, (3) gives
$\rho_{N^{\perp_M}}(A)=|A|-k+\rho_N(Y^\perp)$; hence $A$ is independent
for this rank function exactly when $\rho_N(Y^\perp)=k$, which is the
displayed family.  A basis has size $r-k$ and is precisely the ordinary
support obtained from (2).

If $N_p$ is a quotient of $N_q$, rank increments in $N_p$ are at most
those in $N_q$.  For flats $X\le Y$, applying (3) to
$Y^\perp\le X^\perp$ gives
\[
\rho_{N_q^{\perp_M}}(Y)-\rho_{N_q^{\perp_M}}(X)
 \le
\rho_{N_p^{\perp_M}}(Y)-\rho_{N_p^{\perp_M}}(X).
\]
This is the strong-map criterion; arbitrary subsets reduce to
$M$-closures.
\end{proof}

\begin{lemma}[elementary opposition]\label{lem:elementary}
Let $p\preceq q\preceq\mu_M$, with $\rk p=d$ and $\rk q=d+1$.  Then
$q^{\perp_M}$ and $p^{\perp_M}$ satisfy all adjacent local three-term
relations of ranks $r-d-1$ and $r-d$.
\end{lemma}

\begin{proof}
Fix $A$ of size $r-d-2$ and distinct $x,y,z\notin A$; if $d=r-1$ there
is nothing to check.  Put $U=\cl(A)$, $W=\cl(Axyz)$, and
$\delta=\rk W-\rk U$.  If $A$ is dependent, or if $\delta\le1$, all
three transformed terms are $\infty$.

Assume $\delta=2$.  Let $G=W^\perp$ and $H=U^\perp$, of ranks $d$ and
$d+2$.  A non-$\infty$ term has common $p$-coordinate $\bar p(G)$ and a
$q$-coordinate $\bar q((\cl(Ae))^\perp)$.  In the rank-two interval
$[U,W]$, if one of $x,y,z$ is a loop over $U$, then every transformed
term involving it is $\infty$; if exactly two of the remaining points
coincide, the two corresponding finite terms have the same
$q$-coordinate.  In the only remaining case the three points are
distinct.  Choose a basis $S$ of $G$, representatives
$\ell_x,\ell_y,\ell_z$ in the three dual points of $[G,H]$, and a
complement $C$ extending $H$ to an $M$-basis.  The nested relation for
$q\preceq\mu_M$ with lower set $S$ and upper set
$S\cup\{\ell_x,\ell_y,\ell_z\}\cup C$ has exactly these three finite
$q$-terms.  Adding $\bar p(G)$ gives the desired relation.

Assume $\delta=3$.  Then $x,y,z$ form a rank-three frame over $U$.  Let
$G=W^\perp$ and set
\[
 P_x=(\cl(Ayz))^\perp,\quad P_y=(\cl(Axz))^\perp,\quad
 P_z=(\cl(Axy))^\perp .
\]
Modularity gives that $P_x,P_y,P_z$ are the three atoms over $G$ in the
dual frame and
\[
 (\cl(Ax))^\perp=P_y\vee P_z,\quad
 (\cl(Ay))^\perp=P_x\vee P_z,\quad
 (\cl(Az))^\perp=P_x\vee P_y .
\]
Choose a basis $S$ of $G$ and representatives
$a_e\in P_e\setminus G$.  The adjacent nested relation for
$p\preceq q$ applied to $S$ and $a_x,a_y,a_z$ is exactly the
transformed three-term relation.
\end{proof}

\begin{proposition}\label{prop:set}
The formula $\nu\mapsto\nu^{\perp_M}$ defines an involution of the set
underlying $\Dr(M)$.
\end{proposition}

\begin{proof}
Insert $\nu$ into a complete flag
$h_0\preceq h_1\preceq\cdots\preceq h_r=\mu_M$ using (1).  The ordinary
supports of the $h_j$ form a complete flag from the support of $\nu$ to
$M$; by Lemma \ref{lem:ordinary}, their polar supports form the reversed
complete ordinary flag.  Lemma \ref{lem:elementary} applied to each
adjacent pair shows that
\[
        h_r^{\perp_M},h_{r-1}^{\perp_M},\ldots,h_0^{\perp_M}
\]
satisfies all adjacent local three-term relations.  Its endpoints are
the rank $0$ trivial valuation and $\mu_M$.  By Lemma
\ref{lem:localflag} it is a complete valuated flag, so every layer, in
particular $\nu^{\perp_M}$, is a quotient of $\mu_M$.  Involutivity and
projective well-definedness are immediate from (2).
\end{proof}

\section*{3. Full order reversal}

We need one elementary consequence of being a quotient of $\mu_M$.

\begin{lemma}[circuit minima]\label{lem:circuit}
Let $w\preceq\mu_M$ have rank $k\ge1$, let $A$ be an independent
$(k-1)$-set of $M$, and let $C$ be a circuit of $M/A$ with
$|C|\ge2$.  Then
\[
        \min_{i\in C} w(A\cup i)
\]
is attained at least twice.
\end{lemma}

\begin{proof}
For each $i\in C$, the set $A\cup(C-i)$ is independent and spans the
same flat.  Choose $D$ so that $A\cup(C-i)\cup D$ is an $M$-basis.
Apply the nested relation for $w\preceq\mu_M$ with lower set $A$ and
upper set $A\cup C\cup D$.  The finite $\mu_M$-terms are exactly the
terms indexed by $i\in C$, so the displayed minimum is attained at
least twice.
\end{proof}

\begin{theorem}\label{thm:order}
If $p\preceq q\preceq\mu_M$, then
\[
        q^{\perp_M}\preceq p^{\perp_M}.
\]
Consequently (2) induces an order-reversing involution of $\Dr(M)$.
\end{theorem}

\begin{proof}
Let $\rk p=a$ and $\rk q=b$.  Proposition \ref{prop:set} gives that
the transforms are valuated matroids, and Lemma \ref{lem:ordinary}
puts their ordinary supports in the required quotient order.  If
$a=0$, then $p^{\perp_M}=\mu_M$ and the claim is Proposition
\ref{prop:set}; if $b=r$, the lower transformed rank is $0$ and there
are no incidence relations.  Hence assume $0<a\le b<r$.  By Lemma
\ref{lem:localflag}, it remains to check nested incidence for
$q^{\perp_M}\preceq p^{\perp_M}$.

Fix $A\subseteq T$ with
\[
        |A|=r-b-1,\qquad |T|=r-a+1,
\]
and set $B=T\setminus A$, $m=|B|=b-a+2$.  If $A$ is dependent, all
terms are $\infty$.  Otherwise put $U=\cl(A)$ and $W=\cl(T)$.  A term
can be finite only if $T-i$ is independent; hence, if
$\delta=\rk W-\rk U\le m-2$, all terms are $\infty$.

Suppose first that $\delta=m$.  Then $T$ is independent.  Let
$G=W^\perp$ and $H=U^\perp$, of ranks $a-1$ and $b+1$.  For
$i\in B$ define
\[
        X_i=(\cl(T-i))^\perp,\qquad
        Y_i=(\cl(Ai))^\perp .
\]
Since $B$ is independent over $U$, the flats $\cl(A\cup J)$,
$J\subseteq B$, have ranks $\rk U+|J|$ and form the Boolean sublattice
of $[U,W]$ generated by the frame $B$.  After applying $^\perp$, the
$X_i$ are the atoms over $G$ in the dual frame and
$Y_i=\bigvee_{j\ne i}X_j$.  Choose a basis $S$ of $G$ and
representatives $x_i\in X_i\setminus G$.  Then
$S\cup\{x_i:i\in B\}$ is independent of size $b+1$, and the nested
relation for $p\preceq q$ with this upper set has terms
\[
        p(Sx_i)+q\bigl(S\cup\{x_j:j\ne i\}\bigr)
        =\bar p(X_i)+\bar q(Y_i),
\]
which are exactly
$p^{\perp_M}(T-i)+q^{\perp_M}(Ai)$.  Hence the transformed minimum is
attained at least twice.

It remains to consider $\delta=m-1$.  In the contraction $M/A$, the
set $B$ has nullity one, hence has a unique circuit $C$.  If $C$ is a
singleton loop, every term is $\infty$.  Otherwise $|C|\ge2$.  For
$i\notin C$, the set $T-i$ is dependent, so the transformed term is
$\infty$.  For $i\in C$, the set $T-i$ is independent and has closure
$W$, so the $p^{\perp_M}(T-i)$ factor is the common value
$\bar p(W^\perp)$.  Since $q^{\perp_M}\preceq\mu_M$, Lemma
\ref{lem:circuit} applied to $w=q^{\perp_M}$ and to the circuit
$C$ of $M/A$ shows that the remaining $q^{\perp_M}(Ai)$ minimum is
attained at least twice.  Adding the common summand proves the nested
relation.
\end{proof}

\section*{4. Agreement with the flat involution}

Let $F$ be a flat of rank $f$, and let $q_F$ be the trivial rank
$r-f$ valuated matroid of $M/F\oplus U_{0,F}$.  For an independent
$(r-f)$-set $B$ with $Z=\cl(B)$,
\[
        q_F(B)<\infty\quad\Longleftrightarrow\quad Z\vee F=\hat1.
                                                               \tag{4}
\]
For an independent $f$-set $A$ with $Y=\cl(A)$, (2) and (4) give
\[
 q_F^{\perp_M}(A)<\infty
 \quad\Longleftrightarrow\quad
        Y^\perp\vee F=\hat 1 .
\]
Since $\rk(Y^\perp)=r-f$ and $\rk F=f$, modularity makes this
equivalent to $Y^\perp\wedge F=\hat0$.  Applying the order-reversing
involution gives $Y\vee F^\perp=\hat1$, which is precisely the basis
condition for $M/F^\perp\oplus U_{0,F^\perp}$.  The finite values are
zero on both sides.  Therefore
\[
 \Phi\bigl(\Trop(M/F\oplus U_{0,F})\bigr)
 =
 \Trop(M/F^\perp\oplus U_{0,F^\perp}).
\]

Combining this with Theorem \ref{thm:order}, the answer is yes.  The
desired order-reversing involution is represented on valuated Plucker
coordinates by
\[
        \boxed{\ \nu\longmapsto \nu^{\perp_M}\ } .
\]

\begin{thebibliography}{10}

\bibitem{BEZ}
M.~Brandt, C.~Eur and L.~Zhang,
\emph{Tropical flag varieties},
Advances in Mathematics \textbf{384} (2021), 107695.

\bibitem{DW}
A.~Dress and W.~Wenzel,
\emph{Valuated matroids},
Advances in Mathematics \textbf{93} (1992), 214--250.

\bibitem{DT}
A.~Dress and W.~Terhalle,
\emph{Well-layered maps---A class of greedily optimizable set functions},
Applied Mathematics Letters \textbf{8} (1995), 77--80.

\bibitem{JL}
M.~Jarra and O.~Lorscheid,
\emph{Flag matroids with coefficients},
Advances in Mathematics \textbf{436} (2024), 109396.

\bibitem{Murota}
K.~Murota,
\emph{Matroid valuation on independent sets},
Journal of Combinatorial Theory, Series B \textbf{69} (1997), 59--78.

\bibitem{MurotaBook}
K.~Murota,
\emph{Discrete Convex Analysis},
SIAM, Philadelphia, 2003.

\bibitem{MS}
K.~Murota and A.~Shioura,
\emph{$M$-convex function on generalized polymatroid},
Mathematics of Operations Research \textbf{24} (1999), 95--105.

\bibitem{MSsimple}
K.~Murota and A.~Shioura,
\emph{Simpler exchange axioms for $M$-concave functions on generalized
polymatroids},
Japan Journal of Industrial and Applied Mathematics \textbf{35} (2018),
235--259.

\bibitem{Speyer}
D.~Speyer,
\emph{Tropical linear spaces},
SIAM Journal on Discrete Mathematics \textbf{22} (2008), 1527--1558.

\end{thebibliography}

\end{document}
